\begin{solapartat}
\punts{0,1} Nombre màssic: $A = 13$.

\punts{0,1} Nombre atòmic: $Z = 7$.

\punts{0,2} Nombre de neutrons: $N = A - Z = 13 - 7 = 6$.

\punts{0,45} $^{13}_{7}\mathrm{N} \to {}^{13}_{6}\mathrm{C} + {}^{0}_{+1}\mathrm{e} + \nu$

En aquest apartat es penalitzarà l'omissió del neutrí amb 0,15 p.

\destaca{Alternativament} es pot escriure ${}^{0}_{1}e^+$, ${}^{0}_{+1}\beta$ o
${}^{0}_{1}\beta^+$ en lloc de ${}^{0}_{+1}e$ i ${}^{0}_{0}\nu$ en lloc de $\nu$.

\punts{0,4} La reacció ${}^{1}_{1}\mathrm{p}^+ \to {}^{1}_{0}\mathrm{n} + {}^{0}_{+1}\mathrm{e}^+ + {}^{0}_{0}\nu$
no pot tenir lloc fora d'un nucli ja que la massa del neutró és més gran que la del protó,
i, per tant, cal aportar energia perquè el protó es transformi en un neutró. Si el sistema
està aïllat, llavors no es pot obtenir l'energia que cal perquè la reacció sigui possible.
\end{solapartat}

\begin{solapartat}
$N(t) = N_0 e^{-\lambda t}$

\punts{0,25} $\dfrac{N_0}{2} = N_0 e^{-\lambda T_{1/2}} \Rightarrow \dfrac{1}{2} = e^{-\lambda T_{1/2}} \Rightarrow T_{1/2} = -\dfrac{\ln 0,5}{\lambda} = \dfrac{\ln 2}{\lambda}$

\punts{0,25} $T_{1/2} = \dfrac{\ln 2}{\lambda} \Rightarrow \lambda = \dfrac{\ln 2}{T_{1/2}} = 0,0696\ \mathrm{min^{-1}} = 4,17\ \mathrm{s^{-1}}$ (es pot resoldre el problema
amb unitats de temps de minuts o de segons).

\punts{0,25} $m(t) = m_0 e^{-\lambda t} = 5\,e^{-\lambda t} = 0,620$ ng

\punts{0,25} $A(t) = A_0 e^{-\lambda t}$

\punts{0,25} $\dfrac{A(t)}{A_0} = 0,01 = e^{-\lambda t} \Rightarrow t = -\dfrac{\ln 0,01}{\lambda} = 66,2\ \mathrm{min} = 3972\ \mathrm{s}$
\end{solapartat}
